\begin{mistake}{2026-05-04}{03}

\begin{problemcontent}
设数列 \(\{a_n\}\) 满足
\[
\lim_{n\to\infty}\frac{a_{n+1}}{a_n}=q,
\]
且 \(|q|<1\)，证明：
\[
\lim_{n\to\infty}a_n=0.
\]
\end{problemcontent}

\begin{solutioncontent}
因为 \(|q|<1\)，取常数
\[
r=\frac{1+|q|}{2},
\]
则
\[
|q|<r<1.
\]

由
\[
\lim_{n\to\infty}\frac{a_{n+1}}{a_n}=q,
\]
对 \(\varepsilon=r-|q|>0\)，存在正整数 \(N\)，使得当 \(n\ge N\) 时，
\[
\left|\frac{a_{n+1}}{a_n}-q\right|<r-|q|.
\]
于是当 \(n\ge N\) 时，由三角不等式得
\[
\left|\frac{a_{n+1}}{a_n}\right|
=\left|\frac{a_{n+1}}{a_n}-q+q\right|
\le \left|\frac{a_{n+1}}{a_n}-q\right|+|q|
<r.
\]
因此
\[
|a_{n+1}|<r|a_n|,\qquad n\ge N.
\]

递推可得，对任意 \(k\ge 1\)，
\[
|a_{N+k}|<r^k|a_N|.
\]
令 \(n=N+k\)，则当 \(n>N\) 时，
\[
|a_n|<|a_N|r^{n-N}.
\]

由于 \(0<r<1\)，有
\[
\lim_{n\to\infty}r^{n-N}=0.
\]
又
\[
0\le |a_n|<|a_N|r^{n-N},
\]
由夹逼准则得
\[
\lim_{n\to\infty}|a_n|=0.
\]
因此
\[
\lim_{n\to\infty}a_n=0.
\]
\end{solutioncontent}

\mistakeknowledge{
\begin{itemize}
  \item \textbf{核心公式/定理}：若 \(\lim\limits_{n\to\infty} a_{n+1}/a_n=q\) 且 \(|q|<1\)，则可取 \(|q|<r<1\)，并由极限定义推出充分大时 \(|a_{n+1}/a_n|<r\)。
  \item \textbf{易错点}：不能直接把极限值 \(|q|<1\) 当成每一项比值的上界；必须先取中间常数 \(r\)，再用极限定义得到“从某项开始”的不等式控制。
  \item \textbf{关联知识}：若从某项开始 \(|a_{n+1}|<r|a_n|\)，其中 \(0<r<1\)，则递推得 \(|a_{N+k}|<r^k|a_N|\)，再由等比数列趋零和夹逼准则推出 \(a_n\to0\)。
\end{itemize}}

\end{mistake}
